\documentclass[11pt]{amsart}
\usepackage[T1]{fontenc}
\usepackage{lmodern}
\usepackage{amsmath,amssymb,amsthm}
\usepackage[margin=1.1in]{geometry}
\usepackage{booktabs}
\usepackage{enumitem}
\usepackage{hyperref}

\newtheorem{theorem}{Theorem}[section]
\newtheorem{lemma}[theorem]{Lemma}
\newtheorem{proposition}[theorem]{Proposition}
\newtheorem{corollary}[theorem]{Corollary}
\newtheorem{conjecture}[theorem]{Conjecture}
\theoremstyle{definition}
\newtheorem{definition}[theorem]{Definition}
\newtheorem{remark}[theorem]{Remark}
\newtheorem{example}[theorem]{Example}

\newcommand{\cD}{\mathcal D}
\newcommand{\lean}[1]{\quad{\footnotesize[Lean: \texttt{#1}]}}
\newcommand{\heads}{\operatorname{heads}}
\emergencystretch=3em

\title[Two pairs in Lov\'asz decompositions]{Two pairs in Lov\'asz decompositions:\\
a multi-tail leaf theorem, pre-committed end paths,\\ and the state of $L_N(2)$}
\author{Gallai project (cloud session \texttt{claude/ln})}
\date{28 November 2026 (working draft)}

\begin{document}
\begin{abstract}
Let $F$ be a graph in which every degree is odd. A \emph{Lov\'asz decomposition} of $F$ partitions
$E(F)$ into paths so that every vertex is an end of exactly one path; $Q(v)$ is the path ending at
$v$. Conjecture $L_N$ of route descent asks, for disjoint non-adjacent pairs $c_1a_1,\dots,c_ma_m$,
for a Lov\'asz decomposition with $a_i\notin V(Q(c_i))$ for every $i$. For $m=1$ it is a theorem
($L_1$). The case $m=2$ is the cleanest open statement on the main line of the project
(synthesis~\S32). We do not prove $L_N(2)$. We prove the following, and locate what is left.
\begin{itemize}
\item \emph{Theorem A (multi-tail leaf theorem; Lean).} In a graph with at most one vertex of
positive even degree, any family of \emph{leaf tails} $(c,A)$ (a leaf $c$, and vertices $A$ at
distance $\ge3$ from $c$) can be kept simultaneously, provided no vertex has two neighbours among the
avoided vertices. The leaves need not be sparse, several tails may share a stem, and the even
vertex is unrestricted. This contains the leaf-tail lemma of~\cite{TRN}, the multi-token fin
of~\cite{MT}, and the case of sparse separation~\cite{MS} in which every separated pair is a leaf
and a vertex at distance $\ge3$ from it.
\item \emph{Corollary B (Lean).} $L_N(m)$ holds when all tails but one are leaves at distance
$\ge3$ from their heads and the heads are sparse.
\item \emph{Lemma R (pre-committed end path).} Fix one end path in advance, a path $P$ from $c_i$
avoiding $a_i$. If $P$ ends at $c_j$ and avoids $a_j$, any parity decomposition of $F-E(P)$ finishes
the instance; this is a two-even-vertex existence problem with no constraint at all. If $P$ does
not end at $c_j$, a parity decomposition of $F-E(P)$ with a leaf hung at $c_j$ that keeps a single
leaf tail finishes it. \emph{Criterion R} asks for such a $P$, not ending at $a_j$, together with a
token schedule of~\cite{MT} for what is left; by the transfer theorem of~\cite{MT} it implies
$L_N(2)$ for that instance.
\item \emph{Computation.} Criterion~R holds at all $43{,}400$ instances of order $8$ and at a random
sample of $3{,}000$ instances of order $10$. Every certificate was executed move by move, from an
exact solution of the bottom graph, and the resulting decomposition was checked. By comparison, the two-tail token schedules
of~\cite{MT} cover $72.7\%$ of the instances of order~$8$.
\end{itemize}
We explain, informally, why the natural Fan-type inductions for $L_N(2)$ carry two tokens and two
constraints, and why they stall on dense graphs (the obstruction of~\cite{MT} is a theorem for
specific invariants only). A pre-committed end path removes a constraint before any token moves;
this is what the schedules lacked. The residual open statement is Conjecture~R: every $L_N(2)$
instance satisfies Criterion~R.
\end{abstract}
\maketitle

\section{Introduction}\label{sec:intro}

\begin{conjecture}[$L_N$, route descent~\cite{RD}]\label{conj:LN}
Let $F$ have all degrees odd, and let $c_1a_1,\dots,c_ma_m$ be pairwise disjoint pairs of
non-adjacent vertices. Then $F$ has a Lov\'asz decomposition with $a_i\notin V(Q(c_i))$ for every
$i$.
\end{conjecture}

\emph{Why it matters} (synthesis \S\S31--32). $L_N$ implies Gallai's conjecture for every graph whose
even vertices induce a subgraph with a perfect matching, with every orientation realisable~\cite{RD}.
Through the all-odd matching lifts it gives both spine tests (T1) and (T2) at every root where
$J[N(z)]$ has a perfect matching (\texttt{test1At\_of\_matchingOdd},
\texttt{test2At\_of\_forestBase2}). CLT-B implies it~\cite{CLT}. It holds exhaustively at order
$10$, for two and three fixed pairs (test2, E34).

\emph{What was known.} $L_1$ is a theorem: hang a leaf $\ell$ at $c$, which makes $c$ the only even
vertex, and separate $\ell$ from $a$ by the separation theorem~\cite{RD}. For $m\ge2$ the leaf trick
makes $m$ even vertices. The two-token and $k$-token steering of~\cite{TRN,MT} moves these tokens by
Fan steps until at most one is left, and it is weak on dense instances; an obstruction invariant
explains why (\texttt{not\_steerK}).

\subsection*{Results and status}
\begin{center}\footnotesize
\begin{tabular}{@{}p{0.43\textwidth}p{0.24\textwidth}p{0.27\textwidth}@{}}
\toprule
statement & status & where\\
\midrule
Theorem A (multi-tail leaf theorem) & proved, Lean & \S\ref{sec:A}, \texttt{leaf\_tails}\\
Corollary B (one general tail) & proved, Lean & \S\ref{sec:B}, \texttt{lovasz\_avoid\_leafTails}\\
Lemma R (pre-committed end path) & proved on paper & \S\ref{sec:R}\\
Criterion R implies $L_N(2)$ for the instance & proved on paper (via \cite[Thm.~2.7]{MT}) & \S\ref{sec:R}, Prop.~\ref{prop:R}\\
Criterion R at all order-$8$ instances & computed, every certificate executed & \S\ref{sec:comp}\\
Criterion R at $3{,}000$ random order-$10$ instances & computed, executed & \S\ref{sec:comp}\\
two leaf tails, $\le1$ even vertex, heads not sparse (MTL2) & conjecture; exhaustive $n\le9$ & \S\ref{sec:open}\\
Conjecture R (every instance satisfies Criterion R) & open & \S\ref{sec:open}\\
$L_N(2)$ & \textbf{open} & \\
\bottomrule
\end{tabular}
\end{center}

\section{Preliminaries}\label{sec:prelim}

All graphs are finite and simple. For a path decomposition $\cD$, $b_\cD(v)$ is the number of path
ends at $v$. A \emph{parity decomposition} of $K$ has $b(v)=d_K(v)\bmod 2$ for every $v$. If $K$ has
at most one vertex of positive even degree, a parity decomposition is a Lov\'asz decomposition, and
one exists (Lov\'asz; Lean \texttt{lovasz\_ends}). An \emph{even vertex} has even positive degree; the
even vertices of the current graph are called \emph{tokens}.

\begin{lemma}[Fan's lemma with vertex sets~\cite{Fan05},{\cite[Lemma~2.2]{RD}}]\label{lem:fan}
Let $\cD'$ be a path decomposition of $K'=K-tu$ with $b_{\cD'}(t)$ larger than the number of
neighbours of $u$ in $K'$ with no end. Then $K$ has a path decomposition $\cD$ with
$b_\cD=b_{\cD'}-\mathbf 1_t+\mathbf 1_u$, every path of which lies in $V(P)\cup\{u\}$ for some
$P\in\cD'$.\lean{fan\_single\_supp}
\end{lemma}

A vertex $c$ and a set $A$ are \emph{far apart} if every $a\in A$ is at distance at least three from
$c$: $a\ne c$, $a\not\sim c$, and $a$ and $c$ have no common neighbour.

\begin{definition}[leaf tails]\label{def:tails}
A \emph{leaf tail} of $K$ is a pair $(c,A)$ in which $c$ has degree one and $A$ is far apart from
$c$. For a family $T$ of leaf tails, $\heads T=\bigcup_{(c,A)\in T}A$, and the \emph{tail leaves}
are the vertices $c$. A decomposition \emph{keeps} $T$ if, for every $(c,A)\in T$, no path through
$c$ meets $A$. Since $c$ is a leaf, the path through $c$ is the path ending at $c$. A set $X$ is
\emph{sparse} in $K$ if no vertex of $K$ has two distinct neighbours in $X$.
\lean{LeafTails, KeepsTails, SparseIn}
\end{definition}

Keeping a tail is a property of vertex sets. By Lemma~\ref{lem:fan}, it survives every Fan step
whose target $u$ is neither a head nor a tail leaf: a path through $c$ after the step lies in an old
path through $c$ plus $u$ (\texttt{keepsTails\_of\_cover}).

\section{Theorem A: the multi-tail leaf theorem}\label{sec:A}

\begin{theorem}[multi-tail leaf theorem]\label{thm:A}
Let $K$ have at most one vertex of positive even degree, and let $T$ be a finite family of leaf
tails of $K$ such that $\heads T$ is sparse in $K$. Then $K$ has a Lov\'asz decomposition that keeps
$T$.\lean{leaf\_tails}
\end{theorem}

Nothing is assumed about the tail leaves. Several tails may share a stem, a stem may be adjacent to
the heads of other tails, the even vertex may be a stem or a head, and a head may be the leaf of
another tail. This is what separates Theorem~\ref{thm:A} from three known results.
\begin{itemize}
\item The leaf-tail lemma~\cite[Lemma~6.1]{TRN} (\texttt{avoid\_leafTail}) is the case of one tail
with $|A|=1$ and $A$ odd.
\item The multi-token fin~\cite[Lemma~2.6]{MT} (\texttt{tailGood\_fin}) allows at most one live tail.
\item Sparse separation~\cite{MS} (\texttt{separation\_sparse}), for pairs $\{c,a\}$ with $c$ a
leaf, needs the tail leaves and the heads \emph{together} to be sparse; this fails as soon as a stem
is adjacent to another tail's head, or two tails share a stem. (For larger separated sets sparse
separation says more: it also separates the heads from each other.)
\end{itemize}

\begin{proof}
Induction on $\nu(K)=\sum_v d(v)+[K\text{ has no even vertex}]$ (\texttt{lt\_strong}). Write
$X=\heads T$ and $L$ for the set of tail leaves. Every case below either returns a decomposition
directly or calls the induction hypothesis on a graph $K'$ with $\nu(K')<\nu(K)$, at most one even
vertex, a family $T'$ of leaf tails and $\heads T'\subseteq X$ sparse in $K'$. These conditions pass
to subgraphs, and tails whose leaf becomes isolated are removed from $T'$ (their constraint is
then checked by hand). We check only the new points.

\emph{Case 1: an even vertex $z$ with a neighbour $u\notin X\cup L$} (\texttt{lt\_case\_target}).
Let $K'=K-zu$. In $K'$ only $u$ can be even, $z$ has one end, and every neighbour of $u$ is odd, so
Lemma~\ref{lem:fan} applies with $t=z$. Its target $u$ is neither a head nor a tail leaf, so the
new decomposition keeps $T$.

\emph{Case 2: an even vertex $z$ all of whose neighbours lie in $X\cup L$.} Tail leaves have degree
one, and $z$ has at most one neighbour in $X$, by sparseness. As $d(z)\ge2$, one of the following
holds (the derivation is in \texttt{lt\_strong}):
\begin{itemize}
\item $z$ has two neighbours $c,y$ of degree one; or
\item $N(z)=\{c,y\}$ with $d(c)=1$ (so $d(z)=2$).
\end{itemize}
In both, delete $zc$ and $zy$ (\texttt{lt\_case\_path}). In $K'$ the vertex $c$ is isolated, $y$ is
isolated or even, $z$ is isolated or even, and the alternatives make sure that only one of $y,z$ is
even. Put back the path $c\,z\,y$ (\texttt{PathDecomp.add\_path2\_supp}); the end counts become
those of $K$. The new path keeps every tail: its tail leaves are among $c,y$; for the tail at $c$,
the vertex $z$ is a neighbour of $c$ and $y$ is at distance two, so neither is a head of that tail;
the same holds for a tail at $y$. Every other path is a path of the child.

\emph{Case 3: no even vertex, and a vertex $v$ of degree one} (\texttt{lt\_case\_leaf}). Delete the
edge $vw$ and restore it as a one-edge path; it keeps every tail because a head of a tail at $v$ or
at $w$ is not adjacent to it.

\emph{Case 4: no even vertex, and a vertex $x\notin X$ of degree at least three}
(\texttt{lt\_case\_big}). This is Lov\'asz's Case~II: pick $u\in N(x)$, move the edge $ux$ to a new
leaf at $u$, which makes $x$ the only even vertex; the tails are carried to the new graph, and a tail
leaf equal to $u$ stays a leaf. Delete the new leaf afterwards; then $u$ has $0$ or $2$ ends. With
$0$ ends, add $ux$ as a path, which keeps the tails because $x$ is a neighbour of $u$. With $2$ ends,
use Lemma~\ref{lem:fan} with $t=u$ and target $x$: every neighbour $w\ne u$ of $x$ is odd and has an
end, so the count is $0<2=b(u)$; and $x$ is neither a head nor a tail leaf.

\emph{Case 5: otherwise} (\texttt{lt\_case\_empty}). There is no even vertex and no vertex of degree
one. Every tail leaf has degree one, so $T=\emptyset$ and $X=\emptyset$; as Case~4 does not apply,
every vertex has degree $0$. So there are no edges. (The Lean proof argues with sparseness instead.)
\end{proof}

\begin{remark}[sparseness is used once]\label{rem:sparse}
Sparseness is needed only in Case~2: it caps the number of neighbours of the even vertex in $X$ at
one. Without it the residual is an even vertex $g$ whose neighbours are all heads, except possibly
one tail leaf; for example
$N(g)=\{x_1,x_2\}$ with $x_i$ avoided by tail $i$. With $x_1\not\sim x_2$ one can suppress $g$; the
configurations that resist are the triangle $g\,x_1\,x_2$ and the case where each $x_i$ is the stem
of the other tail. Both reduce to a symmetric form of $L_1$ at distance two, which is open
(the remark on the limits of the leaf trick in~\cite{RD}). Without sparseness Theorem~\ref{thm:A}
is false in general: Section~\ref{sec:comp} gives counterexamples of orders $7$ and $9$ with one even
vertex. For two tails with one head each no counterexample is known (Conjecture~\ref{conj:mtl}).
\end{remark}

\section{Corollary B}\label{sec:B}

\begin{corollary}[one general tail]\label{cor:B}
Let $F$ have no vertex of positive even degree, let $c$ be odd and $a\ne c$ with $a\not\sim c$, and
let $T$ be a family of leaf tails of $F$ with leaves different from $c$. If $\{a\}\cup\heads T$ is
sparse in $F$, then some Lov\'asz decomposition has $a\notin V(Q(c))$ and keeps $T$.
\lean{lovasz\_avoid\_leafTails}
\end{corollary}

\begin{proof}
Hang a leaf $\ell$ at $c$. In $F+c\ell$ the vertex $c$ is the only even vertex, $(\ell,\{a\})$ is a
leaf tail, and the tails of $T$ stay leaf tails; the heads stay sparse because $\ell$ is adjacent
only to $c$. Apply Theorem~\ref{thm:A} and delete $\ell$: the path through $\ell$ becomes $Q(c)$
(the pendant-leaf lemma,~\cite[Lemma~2.3]{NA}).
\end{proof}

So $L_N(m)$ holds whenever all tails but one are leaves of $F$ at distance $\ge3$ from their heads,
and the heads are sparse. With $T=\emptyset$ this is $L_1$.

\section{Why two pairs are a two-token problem}\label{sec:two}

This section is an informal explanation; its only theorems are those cited. A Fan step from a token
$t$ with one end needs the target to have no neighbour without an end (Lemma~\ref{lem:fan}), so it
fails as soon as the target is adjacent to a second token. The schedules of~\cite{MT} carry several
tokens under exactly this restriction. In $L_N(2)$ there are two constraints, and both standard
encodings make tokens or restrict targets.
\begin{itemize}
\item As a leaf tail, the constraint at $c$ needs a leaf at $c$, which makes $c$ a token.
\item As a general tail $\mathrm{av}(c,a)$, it is preserved by a Fan step whose target avoids
$N[c]$ (``Fan transfers carry systems'' in~\cite{CLT}; a sufficient condition).
\end{itemize}
In the schedules we consider, the moving token must therefore avoid the other tokens and the
neighbourhood of the other constraint, and it retires only at a vertex all of whose edges it
deletes. In dense graphs few such moves exist, and for specific invariants~\cite{MT} proves that no
schedule exists (\texttt{not\_steerK}, \texttt{not\_steerK\_triangle}). We checked the alternatives
computationally (Section~\ref{sec:comp}).
\begin{itemize}
\item Token schedules with two tails reach at most one token at $72.7\%$ of the $43{,}400$
instances of order $8$.
\item Using Theorem~\ref{thm:A} as the final step instead of the one-tail fin does not change this
number at order $8$, $m=2$: wherever a schedule reaches one token with sparse heads, some schedule
also drops a tail (both counts are $31{,}532$; \texttt{data/coverage\_order8.txt}). At $m=3$ it adds
$6$ of $4{,}000$ sampled instances. So Theorem~\ref{thm:A} is a structural result; as a
fin for schedules it gains almost nothing.
\item Using the unproved two-leaf-tail statement MTL2 (Conjecture~\ref{conj:mtl}) as the final step
raises the coverage only to $75.7\%$.
\end{itemize}
The missing instances are not exotic. They include the star $K_{1,5}$ with two leaf tails at
distance two, and most dense instances.

\section{Pre-committed end paths}\label{sec:R}

The idea is to fix one end path before any token moves. This removes one constraint outright, at the
price of up to two extra tokens. In case (a) below no tail restricts their targets; in case (b)
only a single leaf tail does.

\begin{lemma}[Lemma R]\label{lem:R}
Let $(F;c_1,a_1;c_2,a_2)$ be an instance of $L_N(2)$, let $\{i,j\}=\{1,2\}$, and let $P$ be a path
of $F$ with at least one edge, one end $c_i$, other end $p$, and $a_i\notin V(P)$. Put $K_0=F-E(P)$.
\begin{enumerate}[label=(\alph*)]
\item If $p=c_j$, $a_j\notin V(P)$, and $K_0$ has a parity decomposition, then $F$ has a Lov\'asz
decomposition with $Q(c_1)=Q(c_2)=P$.
\item If $p\ne c_j$, let $K=K_0+c_j\ell$ with a new leaf $\ell$. If $K$ has a parity decomposition in
which no path contains both $\ell$ and $a_j$, then $F$ has a Lov\'asz decomposition with $Q(c_i)=P$
and $a_j\notin V(Q(c_j))$. The same holds with $K=K_0$ and the leaf $c_j$ itself in place of $\ell$,
when $c_j$ is a leaf of $K_0$.
\end{enumerate}
\end{lemma}

\begin{proof}
In $K_0$ the ends $c_i,p$ of $P$ have lost one edge and the interior vertices of $P$ two. So $c_i,p$
have even degree in $K_0$ and every other vertex has odd degree.

(a) A parity decomposition $\cD_0$ of $K_0$ has no end at $c_i,c_j$ and one end at every other
vertex. $\cD_0\cup\{P\}$ is a path decomposition of $F$ with one end at every vertex, and the path
ending at $c_1$ and at $c_2$ is $P$, which avoids $a_1,a_2$.

(b) The vertex $c_j$ is not an end of $P$, so it has odd degree in $K_0$ and even degree in $K$.
Let $\cD_K$ be the given decomposition. The leaf $\ell$ is an end of its path $P_\ell$, and $c_j$
has no end in $\cD_K$, so $P_\ell$ continues past $c_j$. Deleting $\ell$ gives a parity decomposition
$\cD_0$ of $K_0$ in which the path ending at $c_j$ is $P_\ell-\ell$, which avoids $a_j$, and $c_i,p$
have no end. Then $\cD_0\cup\{P\}$ is a Lov\'asz decomposition of $F$ with $Q(c_i)=P$, which avoids
$a_i$, and with the same path ending at $c_j$, which avoids $a_j$. In the leaf variant no leaf is
added and $P_{c_j}$ is the path ending at $c_j$.
\end{proof}

In (a) the remaining problem has no constraint at all: it asks that the graph $K_0$, which has
exactly the two even vertices $c_1,c_2$ (fewer if one of them had degree one), meet the parity bound.
This is the two-even-vertex problem of the project (connections \S\S O--P). Its Conjecture~31 says
that a $2$-connected graph with two even vertices meets the parity bound unless an odd vertex is
universal. So Conjecture~31 implies $L_N(2)$ for every instance in which $c_1,c_2$ are joined by a
path $P$ avoiding $a_1,a_2$ with $F-E(P)$ $2$-connected and without a vertex of degree $n-1$.

A \emph{schedule} below is a schedule in the sense of~\cite[Def.~2.2]{MT}: a sequence of moves M
(token move), D (drop) and S (suppression) after which at most one token and at most one live tail
are left.

\begin{definition}[Criterion R]\label{def:R}
An instance satisfies \emph{Criterion R} if for some $i$ and some $P$ as in Lemma~\ref{lem:R} one of
the following holds.
\begin{enumerate}[label=(\alph*)]
\item $p=c_j$, $a_j\notin V(P)$, and $(K_0,\emptyset)$ has a schedule.
\item $p\notin\{c_j,a_j\}$, and $(K,\{\tau\})$ has a schedule, where either $K=K_0+c_j\ell$ and
$\tau=(\ell,c_j,a_j)$, or $c_j$ is a leaf of $K_0$ with stem $w$, $a_j\notin\{w\}\cup N_{K_0}(w)$,
$K=K_0$ and $\tau=(c_j,w,a_j)$.
\end{enumerate}
\end{definition}

\begin{proposition}\label{prop:R}
An instance that satisfies Criterion R satisfies $L_N(2)$.
\end{proposition}

\begin{proof}
In (a) there is no tail, and the transfer theorem~\cite[Thm.~2.7]{MT} (\texttt{steerK\_good}) gives
a parity decomposition of $K_0$; Lemma~\ref{lem:R}(a) finishes. In (b), recall that a tail
of~\cite[Def.~2.1]{MT} is a triple $(c,s,x)$ of a leaf, its stem and the avoided vertex; the single
tail $\tau$ is live in $K$, as its leaf edge is present. (TS) of~\cite{MT} holds because
$a_j\notin\{\ell,c_j\}$ (respectively $a_j\notin\{c_j,w\}$). (TI) holds in $K$: the leaf has a
single neighbour, $a_j\not\sim c_j$ (respectively $a_j\not\sim w$), and $a_j$ is odd in $K$. For the
last point, $a_j\ne c_j$ (so the leaf edge at $c_j$ does not change its degree), $a_j\ne c_i$
because the pairs are disjoint, and $a_j\ne p$; so $a_j$ is either off $P$ or an interior vertex of
$P$, which loses two edges. The transfer theorem gives a parity
decomposition of $K$ in which no path contains both the leaf of $\tau$ and $a_j$. This is the
hypothesis of Lemma~\ref{lem:R}(b), which finishes.
\end{proof}

The condition $p\ne a_j$ is needed for this argument (not for the truth of $L_N(2)$). If $p=a_j$
and $d_F(a_j)\ge3$, then $a_j$ is a token, (TI) fails, and a schedule can suppress $a_j$ and put it
on the leaf path. Allowing $P$ to end at $a_j$, the independent
mutation test of audit~A40 found $317$ failing executions among $2{,}808$ schedules on random
instances of order $8$; this is also the bug described in Section~\ref{sec:comp}. Our certificates
do not rely on Proposition~\ref{prop:R}: each one is executed move by move, the bottom is solved
exactly, and every level is re-checked (Section~\ref{sec:comp}). In (b) the graph $K$ may have three
tokens ($c_i$, $p$ and $c_j$), but only one constraint.

\section{Computations}\label{sec:comp}

\emph{This section is computational; no proof above depends on it.} The scripts and the data are in
the directories \texttt{scripts} and \texttt{data} of \texttt{manuscripts/ln}.

\subsection*{Tools}
\begin{itemize}
\item \texttt{lnsolve.c}: exact depth-first search for parity decompositions with constraints
(the path ending at $t$ avoids $a$; two vertices on no common path; two vertices joined by one path),
with failure memoisation. Validated against the independent brute force of \texttt{lnlib.py}
(exact cover by all simple paths) on random instances of order $\le7$, and with the joined-ends
constraint of orders $6$ and $8$, with $0$ mismatches; every witness is re-checked
(\texttt{validate.py}, \texttt{validate\_join.py}; \texttt{data/validation.txt}).
\item \texttt{gen\_*.py}: instance generators over \texttt{nauty-geng}.
\item \texttt{sched.py}, \texttt{sched3.py}: token schedules (two and three tails).
\item \texttt{routeR.py}: Criterion R. For each instance it searches $i$ and $P$ (shortest first) and
a schedule, then \emph{executes} the certificate independently of the search: it replays the moves,
solves the bottom exactly with \texttt{lnsolve}, undoes the moves (Fan's forced chain for M, the path
$c s$ for D, subdivision for S), checks at every level that the object is a parity decomposition
keeping the live tail, deletes the leaf, adds $P$, and checks the final Lov\'asz decomposition of $F$
with \texttt{lnlib.check\_witness}.
\end{itemize}

\subsection*{Results}
Instances of $L_N(2)$: connected graphs with all degrees odd, and an unordered pair of disjoint
ordered pairs $(c_1,a_1),(c_2,a_2)$ with $c_i\not\sim a_i$. Order $8$: $43{,}400$ instances
(exhaustive). Order $10$: $3{,}000$ instances (\texttt{sample10.py}, seed $11$). Each draw takes a
graph uniformly from the $32{,}851$ connected all-odd graphs and four distinct vertices uniformly,
and is repeated from scratch when $c_1\sim a_1$ or $c_2\sim a_2$. Graphs are therefore weighted by
their number of admissible quadruples, which under-represents dense graphs (mean edge count $21.4$
in the sample against $22.6$ over the graphs; audit~A40 follow-up).

\begin{center}\footnotesize
\begin{tabular}{@{}lrr@{}}
\toprule
route & order $8$ (of $43{,}400$) & status of the route\\
\midrule
two-tail token schedule, one-tail fin~\cite{MT} & $31{,}532$ ($72.7\%$) & proved~\cite{MT}\\
two-tail token schedule, Theorem~A as fin & $31{,}532$ & proved\\
two-tail token schedule, MTL2 as fin & $32{,}837$ ($75.7\%$) & conditional on Conj.~\ref{conj:mtl}\\
joint decomposition $Q(c_1)=Q(c_2)$ avoiding $a_1,a_2$ (exact) & $40{,}538$ ($93.4\%$) & --\\
Lemma R(a) with a free schedule & $40{,}532$ ($93.4\%$) & proved (Prop.~\ref{prop:R})\\
schedules $\cup$ R(a) & $42{,}904$ ($98.9\%$) & proved\\
\textbf{Criterion R} & $\mathbf{43{,}400}$ \textbf{(all)} & proved (Prop.~\ref{prop:R}); executed\\
\bottomrule
\end{tabular}
\end{center}

At order $10$, Criterion R holds at all $3{,}000$ sampled instances. Every certificate at both
orders was executed and the final decomposition verified (\texttt{data/routeR\_order8.txt},
\texttt{data/routeR\_order10\_sample.txt}, with summaries). The certificates found:
\begin{center}\footnotesize
\begin{tabular}{@{}lrrrr@{}}
\toprule
order & instances & (a) joint & (b) one tail & $P$ a single edge\\
\midrule
$8$ (all) & $43{,}400$ & $14{,}853$ & $28{,}547$ & $34{,}688$\\
$10$ (sample) & $3{,}000$ & $737$ & $2{,}263$ & $2{,}778$\\
\bottomrule
\end{tabular}
\end{center}
The search tries $i=1$ first and short paths first, so these counts describe the certificates found,
not minimal ones. Restricted to a single edge $P$ (with either $i$), Criterion R holds at $38{,}908$
of the $43{,}400$ instances of order $8$ (\texttt{r1count.py}); longer paths are needed for the rest. Schedules have up to $14$ moves at order $8$ and $19$ at order $10$.

\emph{A bug found by the executor.} A first version of the search (\texttt{precommit.py}) let $P$ end
at $a_j$. Then $a_j$ is a token, and a suppression of that token can put $a_j$ on the leaf path. The
executor rejected $287$ of its order-$8$ certificates. The search now forbids $P$ to end at $a_j$
(Definition~\ref{def:R}(b)) and never suppresses a head, as in the invariant (TI) of~\cite{MT}; all
numbers above are from the corrected search, and every certificate was executed. Audit~A40
re-checked every certificate at both orders against Definition~\ref{def:R}, including (TI) after
every move, and re-executed each one with its own code.

\emph{Solver validation.} \texttt{lnsolve} agreed with the brute force on $1{,}500$ random instances of
order $\le7$ ($880$ satisfiable, $620$ not), and on $600$ instances of orders $6$ and $8$ with the
joined-ends constraint ($538$ satisfiable, $62$ not); every witness was re-checked
(\texttt{data/validation.txt}). Audit~A40 cross-checked it against an independent solver on $1{,}800$
further instances, with $0$ disagreements.

\emph{The two-leaf-tail statement} (Conjecture~\ref{conj:mtl}) is exhaustive for $n\le9$: every
connected graph with at most one even vertex, every pair of leaf tails $(c_i,w_i,x_i)$ with
$x_i\notin\{c_i,w_i\}\cup N(w_i)$: $8$, $84$, $385$, $11{,}040$ instances at $n=6,\dots,9$ under the
static conditions of~\cite{MT}, and $40$, $483$, $1{,}192$, $38{,}100$ without them (heads may then
be the leaf or the stem of the other tail, or the even vertex); all satisfiable. With $x_i\sim w_i$ allowed it fails ($12$ instances at $n=7$, all with the stem even
of degree two or a similar local defect), so the distance condition is needed when there is an
even vertex.

\emph{Without sparseness, larger head sets fail} (found by audit~A40 in an exhaustive search over
$n\le9$, and confirmed with \texttt{lnsolve}). Two examples, each with one even vertex:
\begin{itemize}
\item $n=7$, edges $03,05,06,14,15,16,26,35$ (the even vertex is $3$), tails $(2,\{5\})$ and
$(4,\{0,2\})$. Here one tail avoids the other tail's leaf.
\item $n=9$, leaves $4,5,6$ at stems $0,1,2$; each of $0,1,2,3$ is joined to both $7$ and $8$, and
$78$ is an edge (the even vertex is $3$); tails $(4,\{1,2,3\})$, $(5,\{0,2,3\})$, $(6,\{0,1,3\})$. Here
no tail avoids another tail's leaf.
\end{itemize}
An earlier draft reported, from a sample, that such families never failed as long as no tail
avoids another tail's leaf; the second example refutes this. Conjecture MTL2 (two tails, one head
each) survives the same exhaustive search: $15{,}690$ instances at $n\le8$ and $121{,}205$ at $n=9$,
over all graphs including disconnected ones.

\section{What remains}\label{sec:open}

\begin{conjecture}[MTL2]\label{conj:mtl}
Theorem~\ref{thm:A} holds without the sparseness of the heads, for two tails with one head each.
\end{conjecture}

\begin{conjecture}[Conjecture R]\label{conj:R}
Every instance of $L_N(2)$ satisfies Criterion R.
\end{conjecture}

By Proposition~\ref{prop:R}, Conjecture~\ref{conj:R} implies $L_N(2)$. It is a statement about token
schedules with at most one constraint, which is the regime where steering is strong; the open part is
to choose $P$ so that a schedule provably exists. Two partial observations point the way.
\begin{itemize}
\item If $c_i$ has degree one, the token at $c_i$ disappears when $P$ is removed. If $c_i$ has
degree three and its two remaining neighbours are non-adjacent and meet the side conditions of
move S, it can be suppressed (not in a triangle). Then only $p$ and $c_j$ are tokens.
\item In (a), $P$ can be chosen through every vertex of degree $n-1$, which removes the only known
obstruction (O1) of the two-even problem.
\end{itemize}
For general $m$ the same idea pre-commits $m-1$ end paths; the tokens then multiply, and we have not
tested it.

\section{Formalisation}\label{sec:lean}

The file is \texttt{lean/GallaiLean/LeafTails.lean}; it has no \texttt{sorry}, and the full library
builds. For the headline theorems \texttt{leaf\_tails} and \texttt{lovasz\_avoid\_leafTails},
\texttt{\#print axioms} gives only \texttt{propext}, \texttt{Classical.choice} and
\texttt{Quot.sound} (\texttt{data/axioms.txt}; checked again by audit~A40).

\begin{center}\footnotesize
\begin{tabular}{@{}ll@{}}
\toprule
manuscript & Lean\\
\midrule
Definition~\ref{def:tails} & \texttt{LeafTails}, \texttt{heads}, \texttt{KeepsTails}, \texttt{SparseIn}\\
transfer of kept tails & \texttt{keepsTails\_of\_cover}\\
adding a path of length two & \texttt{GPath.ofAdj2}, \texttt{PathDecomp.add\_path2\_supp}\\
Theorem~\ref{thm:A}, Cases 1--3 & \texttt{lt\_case\_target}, \texttt{lt\_case\_path}, \texttt{lt\_case\_leaf}\\
Theorem~\ref{thm:A}, Cases 4--5 & \texttt{lt\_case\_big}, \texttt{lt\_case\_empty}\\
Theorem~\ref{thm:A} & \texttt{lt\_strong}, \texttt{leaf\_tails}\\
Corollary~\ref{cor:B} & \texttt{lovasz\_avoid\_leafTails}, \texttt{KeepsTails.avoid}\\
\bottomrule
\end{tabular}
\end{center}
Lemma~\ref{lem:R} and Proposition~\ref{prop:R} are not formalised. The proof of the lemma uses
only end counting and the pendant-leaf lemma, both of which are in the library
(\texttt{PathDecomp.add\_path} and \texttt{Retract.delete\_supp}). The proposition adds the
transfer theorem, which is also in Lean (\texttt{steerK\_good}).

\begin{thebibliography}{9}
\bibitem{CLT} Gallai project, \emph{Constrained Lov\'asz decompositions}, \texttt{manuscripts/clt}.
\bibitem{Fan05} G.~Fan, Path decompositions and Gallai's conjecture, \emph{J. Combin. Theory Ser. B}
93 (2005), 117--125.
\bibitem{MS} Gallai project, \emph{Multi-pair separation}, \texttt{manuscripts/multi-sep}.
\bibitem{MT} Gallai project, \emph{Multi-token steering}, \texttt{manuscripts/multi-token}.
\bibitem{NA} Gallai project, \emph{Neighbour avoidance}, \texttt{manuscripts/na}.
\bibitem{RD} Gallai project, \emph{Route descent}, \texttt{manuscripts/route-descent}.
\bibitem{TRN} Gallai project, \emph{The tight transfer}, \texttt{manuscripts/trn}.
\end{thebibliography}
\end{document}
